\section{Lesson 6, 05/10}
\subsection{Massive and massless particles representations}
\textbf{Massive particles}\\
It is known from SR that if a particle is massive, there exist a reference frame
in which the particle is at rest, i.e.
$$P^\mu = (m,\mathbf{0})$$
which means that
$$W^0 = 0\quad W^i = -\frac{1}{2} \varepsilon^{0ijk}M_{jk}$$
recalling the definition of $J^i$ this is
$$W^i = -\frac{1}{2}mJ^i$$
since the particle is at rest, this cannot represent an orbital angular momentum,
so it must be intrinsic, i.e. \textit{spin}:
$$W^i = -\frac{1}{2}mS^i$$
and so massive particles are identified by Casimir operators $P^2$ and $S^2\propto
W^2$ and by the commuting operators $S_3$ and $P^\mu$:
there are, for each value of $m,s$, $2s+1$ internal degrees of freedom.

\textbf{Massless particles}\\
If the particle is massless,
$$P_\mu P^\mu = m^2 = 0$$
which means
$$W_\mu W^\mu = 0,$$
and in general
$$P_\mu W^\mu = 0:$$
from these 3 conditions it can be shown\footnote{moodle?} that
$$W^\mu = h P^\mu$$
where $h = \pm 1$ is called \textbf{helicity} and represents the projection of the
angular momentum on the direction of the momentum
$$h = \frac{\mathbf{J}\cdot\mathbf{P}}{|\mathbf{P}|};$$
in this case, they can only be parallel or anti-parallel: these states are called
\textit{polarizations} and, for photons, they correspond to light polarizations.

This framework also allows a sector of the theory in which spin is continuous
($m=0$, $W^2<0$ states are like this, but they have not been observed in nature)
and a sector in which some particles, tachyons can move at $v>c$.

The key takeaway from this arguments is that

\noindent\fbox{%
    \parbox{\textwidth}{%
        \textbf{The irreducible representations of the Poincaré Group are in
        one-to-one correspondence with the possible types of elementary particles.\\
        The mathematical structure of space-time symmetries determines what kinds
        of particles can exist in nature.}
    }%
}
\subsection{Non-finite representations of groups: field representations}
Until this point, all the representations have been finite-dimensional, with
generators represented as matrices, but a field is continuous.

\textbf{Def.} A field is a well-behaved function of spacetime
$$\phi_a(x^\mu) = \begin{pmatrix}
    \phi_1(x^\mu)\\
    \phi_2(x^\mu)\\
    ...
\end{pmatrix}
$$
matrices can transform a field into another with linear combinations of the old
components as new components: these transformations are finite-dimensional, while
the ones on the coordinates $x^\mu$ can also be infinite-dimensional.

In QM the allowed transformations were unitary transformations, as they preserve
squared norms, which are the observable part of the theory: these transformations
are infinite-dimensional (think of the time evolution: its generator is the
Hamiltonian $H$, which is a differential operator).

In parallel to the definition of the angular momentum operator in QM, the Lorentz
generator for infinite-dimensional transformation is
$$M_{\mu\nu} = i (x_\mu\partial_\nu-x_\nu\partial_\mu)$$
and so a field transforms like
$$\phi_a(x) \rightarrow \phi_a'(x') = M\indices{^a_b}(\Lambda)\phi_b(\Lambda x)$$
where $\Lambda$ are infinite-dimensional transformations and $M\indices{^a_b}$ are
finite-dimensional;
the total transformation can be expressed with the exponential map
$$\phi_a(x')=\exp\left[-\frac{i}{2}\omega^{\rho\sigma}M_{\rho\sigma}\right]\phi_b(x),
$$
with
$$M_{\mu\nu} = M_{\mu\nu}^\text{inf} + M_{\mu\nu}^\text{fin}:$$
where the finite part generates the transformations of the field components\footnote{
Q: why are we assuming that the parameters are the same number and are the same?}
and the infinite part those of the space-time coordinates. This generators can be
expressed together thanks to the BCH formula.
